\section{Introduction}

A convex body in $\mathbb R^n$ is a compact convex set with nonempty interior.
A set $X\subset\mathbb R^n$ is \emph{locally finite} if every bounded set
contains only finitely many of its points. It gives a \emph{translative
covering} of $K$ when $K+X=\mathbb R^n$. Write $Q_R=[-R,R]^n$. We define
the translative covering density by
\begin{equation}\label{eq:unrestricted-density}
 \theta_T(K)=\inf_{\substack{X\text{ locally finite}\\K+X=\mathbb R^n}}
 \vol(K)\limsup_{R\to\infty}
       \frac{\#(X\cap Q_R)}{(2R)^n}.
\end{equation}
Thus the density measures the asymptotic volume of translates per unit
volume, counting overlaps. It is not the fraction of space covered, which
is one for every covering. No asymptotic density limit is assumed for an
individual arrangement.

A \emph{periodic translative covering} is a covering
\[
 \mathbb R^n=K+\Lambda+T,
\]
where $\Lambda$ is a full-rank lattice and $T$ is a finite set of distinct
representatives modulo $\Lambda$. Its center intensity is
$|T|/\det\Lambda$, and its density is $|T|\vol(K)/\det\Lambda$. Put
\[
 \theta_{\mathrm{per}}(K)=\inf\left\{
 \frac{|T|\vol(K)}{\det\Lambda}:K+\Lambda+T=\mathbb R^n
 \right\}.
\]
The infimum ranges over every lattice and every finite set $T$. In particular, there is no bound on the number of cosets.
Allowing repeated representatives of the same coset gives the same
infimum: deleting redundant representatives preserves the covering and
does not increase its listed density. When we later retain indexed center
lists, we count such repetitions explicitly.
Both infima are finite: an interior cube of $K$ gives a covering by
translates on a sufficiently fine cubic lattice.

The classical periodization principle gives
\begin{equation}\label{eq:density-equivalence}
 \theta_T(K)=\theta_{\mathrm{per}}(K).
\end{equation}
We prove the precise statement in Appendix~\ref{sec:periodization},
including agreement with lower or upper centered-ball density and average
covering multiplicity. See also Prosanov~\cite[Section~1.1 and
Theorem~2.1]{Prosanov2021PackingCovering}, who attributes the principle to
Rogers. These are equalities of optimal densities; the lower and upper
densities of an individual nonperiodic arrangement need not agree.

\begin{theorem}\label{thm:main}
There are absolute constants $c>0$ and $n_0$ such that, for every integer
$n\ge n_0$, some centrally symmetric convex body $K_n\subset\mathbb R^n$
satisfies
\[
 \theta_T(K_n)>c n\log n.
\]
\end{theorem}

Thus there is no absolute constant $C$ such that
$\theta_T(K)\le Cn$ for every convex body in every dimension, even if the
bodies are required to be centrally symmetric. This answers negatively
the universal linear-bound question recorded by
Nasz\'odi~\cite[Section~3]{Naszodi2018Flavors}.
Rogers~\cite{Rogers1957Coverings} proved that every convex body in dimension
$n\ge3$ has a periodic translative covering of density at most
\begin{equation}\label{eq:rogers}
 n\log n+n\log\log n+5n.
\end{equation}
Together with Theorem~\ref{thm:main}, this shows that the supremum of
covering density over convex bodies has order $n\log n$. The same order holds when
one restricts the bodies to be centrally symmetric. The constant in
Theorem~\ref{thm:main} is not optimized.
The bodies may moreover be chosen to be centrally symmetric polytopes
with rational vertices, with the same constant $c$; see
Corollary~\ref{cor:rational-polytope}.

\subsection{Antecedents and the additional uniformity}

Rogers' theorem supplies a universal upper bound for arbitrary convex
bodies. Fejes T\'oth~\cite{FejesToth2009Covering} sharpened its lower-order
term to obtain density
$n(\log n+\log\log n+1+o(1))$ using only $O(\log n)$ cosets of a common
lattice. A finite union of lattice cosets has substantially more freedom
than one lattice, even with this bound on the number of cosets.
On the lower-bound side, the classical work of Coxeter, Few and
Rogers~\cite{CoxeterFewRogers1959Spheres} gives an order-$n$ obstruction
already for unrestricted coverings by equal Euclidean balls.
For Euclidean balls, Dumer~\cite[Corollary~2]{Dumer2007Spheres} gives the
upper bound $(1/2+o(1))n\log n$. The gap for Euclidean balls remains;
our theorem determines the optimal order when the convex body may also
be chosen.

The proof builds on the random-slab construction and Poisson witness method
of Li and Liu~\cite[Theorem~1.3]{LiLiu2026Lattice}. They prove an order-$n\log n$ lower bound
for \emph{lattice} covering density, in which the centers form a single
lattice. Their common-intersection estimates, terminal clustering of caps,
and overlap-diagram bounds provide the local probabilistic machinery used
here. We give the necessary versions and their proofs, including the
continuous form of Janson's inequality.

The additional issue is to exclude every finite union of lattice cosets.
A net of lattice bases cannot encode an arbitrary number of cosets. We
instead localize an arrangement in a fixed Euclidean window, keeping every
center with its multiplicity, and round the resulting list to a finite grid.
Periodic averaging controls this list even when the lattice has very short
periods or a highly skew fundamental domain. The resulting finite class is
fixed before the body is sampled, so it includes arrangements chosen after
the body is known. A sufficiently strong lower-tail bound then excludes
all of them simultaneously.

For comparison, Bukh, Gao, Liu, Pikhurko, and Sun~\cite[Theorem~1.3]{BukhEtAl2025Random}
show that, in their packing-torus model, Poisson centers at a specified
density of order $n\log n$ fail to cover a given body with probability
tending to one.
Their random-centers barrier has different quantifiers: an exceptional arrangement might still
cover efficiently. Here the body is random and the conclusion excludes
every periodic arrangement for one realization of that body.

\subsection{Proof outline}

Intersect a slightly enlarged Euclidean ball with symmetric slabs whose unit
normals are sampled by a Poisson process.
The normalization gives this body a substantial volume with positive
probability. A hypothetical covering of low center intensity supplies,
after localization and rounding, a finite center list with many suitable
target points in a small window. We must show that each such list is very
unlikely to cover that window.

At a target $y$, only centers $p$ within the outer ball's radius of $y$
can cover it. For each such center, the residual vector $y-p$ determines
the antipodal caps of slab normals that would remove $y$ from that
translate. A hole appears if the Poisson process hits at least one cap
in each pair. Nearly coincident cutting sets are strongly dependent.
Following Li and Liu, we group these conditions into blocks and choose one common subset of normals
for each block. Assigning distinct Poisson points to these subsets certifies
a hole. The block construction controls the complete overlap sum of these
certificates, including their self-overlaps.

The targets are chosen at two angular scales. For every pair of distinct
targets, every pair of their relevant residual directions is separated
at a scale of order $\sqrt{\log n/n}$, identifying opposite directions
because the slabs are symmetric. The cap estimates then bound the
dependence between their hole certificates. The body lies in a ball
whose radius is only slightly larger than the slab half-width, so its
cutting caps have a small fixed angular radius. At a larger fixed
separation, all cutting caps belonging to one target are disjoint from
those belonging to the other. A cylinder-volume estimate lets us arrange
this stronger separation for all but an exponentially small fraction of
target pairs, while retaining the finer separation for every pair.

Let $R$ denote the density lower bound sought. We can choose
$\exp(c_1n\log n)$ targets for a fixed $c_1>0$, enough to absorb a
self-overlap cost $\exp(CR)$ when $R=c n\log n$ and $c$ is sufficiently
small. The two angular scales control the remaining overlap between
distinct targets. Janson's lower-tail inequality turns these bounds into
a covering probability that is doubly exponentially small in $n$, with
decay strong enough to dominate the number of rounded patterns.
A second-moment bound would not give the required decay.

Section~\ref{sec:construction} defines the body and the constants.
Section~\ref{sec:localization} constructs the finite pattern class.
Sections~\ref{sec:caps}--\ref{sec:poisson} establish the cap, block, and
Poisson estimates. Section~\ref{sec:targets} proves the fixed-pattern bound,
and Section~\ref{sec:assembly} completes the proof.
Section~\ref{sec:encoding} gives rational-polytope and rational-normal
versions of the construction.
Appendix~\ref{sec:periodization} proves the density equivalence that
transfers the periodic exclusion to unrestricted coverings.

\subsection{Notation}

Logarithms are natural, and $[t]_+=\max\{t,0\}$. We write $|x|$ for Euclidean norm,
$B_r=\{x\in\mathbb R^n:|x|\le r\}$, $B=B_1$, and
$\omega_n=\vol(B)$. For measurable sets in $\mathbb R^n$, $|A|$ also denotes
Lebesgue volume; for finite lists and index sets it denotes cardinality.
Surface measure $\sigma$ on $S^{n-1}$ is normalized to have total mass one.
For nonzero vectors we use the projective angle
\[
 \angle(x,z)=\arccos\frac{|x\cdot z|}{|x||z|}\in[0,\pi/2].
\]
Lists are indexed: two entries can occupy the same point or define the same
cap. Every sum and cardinality over such a list retains its multiplicities.
Constants implicit in $O(\cdot)$ depend only on parameters explicitly fixed
before $n$ tends to infinity.
